\documentclass[conference]{IEEEtran}

% ── Packages ─────────────────────────────────────────────────────────────────
\usepackage[T1]{fontenc}
\usepackage[utf8]{inputenc}
\usepackage{cite}
\usepackage{amsmath,amssymb}
\usepackage{graphicx}
\usepackage{url}
\usepackage{hyperref}
\usepackage{booktabs}
\usepackage{array}
\usepackage{listings}
\usepackage{xcolor}
\usepackage{float}       % provides [H] placement

% ── Code listing style ────────────────────────────────────────────────────────
\lstset{
  basicstyle=\ttfamily\footnotesize,
  keywordstyle=\color{blue}\bfseries,
  commentstyle=\color{gray},
  breaklines=true,
  frame=single,
  numbers=left,
  numberstyle=\tiny\color{gray},
  language=C
}

% ── Title ─────────────────────────────────────────────────────────────────────
\title{Concurrent Web Server: Thread Pool Architecture\\
with Size-Based Scheduling}

\author{%
  \begin{tabular}{cc}
    \begin{minipage}[t]{0.60\columnwidth}\centering
      \textbf{Ishit Arhatia}\\
      \small MS Data Science\\
      \small Khoury College of Computer Sciences\\
      \small Northeastern University
    \end{minipage}
    &
    \begin{minipage}[t]{0.60\columnwidth}\centering
      \textbf{Rudrajit Banerjee}\\
      \small MS Computer Science\\
      \small Khoury College of Computer Sciences\\
      \small Northeastern University
    \end{minipage}
    \\[2em]
    \begin{minipage}[t]{0.60\columnwidth}\centering
      \textbf{Helena Zhuo}\\
      \small MS Computer Science\\
      \small Khoury College of Computer Sciences\\
      \small Northeastern University
    \end{minipage}
    &
    \begin{minipage}[t]{0.60\columnwidth}\centering
      \textbf{Tomas Colorado}\\
      \small MS Computer Science\\
      \small Khoury College of Computer Sciences\\
      \small Northeastern University
    \end{minipage}
  \end{tabular}
}

\begin{document}

\maketitle

% ═════════════════════════════════════════════════════════════════════════════
\begin{abstract}
% ═════════════════════════════════════════════════════════════════════════════

This paper presents a performance evaluation of size-based request scheduling
in a multi-threaded web server with producer-consumer architecture, built
around a fixed-size thread pool and a bounded buffer, following the
architecture described in the OSTEP curriculum~\cite{ostep}. The main thread
continuously accepts incoming HTTP connections and enqueues requests into a
shared buffer; a pool of worker threads dequeue and service those requests
according to a configurable scheduling policy. Our research extension
investigates \emph{size-based scheduling}---specifically comparing
First-Come-First-Served (FCFS) against Shortest-File-First (SFF)---to evaluate
their impact on mean response time, tail latency, and system behavior under
varying load conditions. Drawing on established research demonstrating that web
file sizes follow heavy-tailed distributions where most requests are small but
most bytes originate from a small number of large files~\cite{harchol2003}, we
hypothesize that SFF scheduling will significantly reduce mean response time
without substantially penalizing large-file requests. We validate this
hypothesis through controlled experiments measuring throughput, response time
distribution (mean, p50, p95, p99), and system behavior across different
levels of concurrent connections and file-size mixes.

\end{abstract}

\begin{IEEEkeywords}
Concurrent web server, thread pool, bounded buffer, producer-consumer,
scheduling policy, First-Come-First-Served (FCFS), Shortest-File-First (SFF),
size-based scheduling, heavy-tailed workloads, POSIX threads.
\end{IEEEkeywords}

% ═════════════════════════════════════════════════════════════════════════════
\section{Introduction}
% ═════════════════════════════════════════════════════════════════════════════

Modern web servers must handle thousands of simultaneous client connections
while maintaining low latency and high throughput. A foundational design choice
in any such system is how to move work from the single point of connection
acceptance---a listening socket---to a pool of threads that can service
requests in parallel. The canonical solution, described by Welsh et al.\ in the
SEDA architecture~\cite{welsh2001} and formalized in the OSTEP curriculum,
is the \emph{producer-consumer bounded buffer}: the main thread produces
connection descriptors into a fixed-capacity queue, and a fixed-size pool of
worker threads consume from it. This decoupling keeps the acceptance loop
fast and provides natural backpressure when workers are saturated.

Once this architecture is in place, a second design choice emerges: in what
\emph{order} should workers dequeue requests? The simplest policy is FCFS
(First-Come-First-Served), which processes requests strictly in arrival order.
FCFS is fair, simple to implement, and free of starvation. However, web
workloads are not uniform. Crovella and Bestavros showed that World Wide Web
traffic is self-similar and that file sizes follow heavy-tailed (Pareto-like)
distributions~\cite{crovella1997}: the vast majority of objects are small
(a few kilobytes), yet a small number of large files account for most bytes
transferred. Under these conditions, FCFS performs poorly---a single large
file can block dozens of waiting small requests behind it.

Size-based scheduling policies such as Shortest-File-First (SFF) exploit this
skew. By allowing workers to select the smallest pending request from the
buffer rather than the oldest, small requests jump ahead of large ones,
completing quickly and clearing queue capacity. Harchol-Balter et al.\
proved theoretically and through trace-driven simulation that this approach
can reduce mean response time by orders of magnitude under moderate-to-high
load without materially starving large requests~\cite{harchol2003}. Schroeder
and Harchol-Balter further showed that the performance gap widens dramatically
as server load approaches capacity~\cite{schroeder2006}---exactly the operating
regime that matters most.

This paper asks: \emph{does SFF scheduling meaningfully outperform FCFS in a
POSIX thread pool web server serving a realistic heavy-tailed workload, and
what are the fairness trade-offs for large-file clients?} We answer this
question by building a fully conformant G.5 OSTEP web server and subjecting it
to automated load tests using an external benchmarking rig.

The remainder of this paper is organized as follows. Section~\ref{sec:related}
surveys the relevant literature grouped by theme. Section~\ref{sec:design}
describes our system architecture and key design decisions.
Section~\ref{sec:methodology} defines our experimental methodology and
benchmarking rig. Section~\ref{sec:results} presents baseline results and
preliminary findings. Section~\ref{sec:conclusion} summarizes our work and
outlines next steps.

% ═════════════════════════════════════════════════════════════════════════════
\section{Related Work}
\label{sec:related}
% ═════════════════════════════════════════════════════════════════════════════

Related work naturally divides into three themes: the architecture of
concurrent servers, the theory of size-based scheduling, and modern
production-scale scheduling systems that validate or extend these ideas.

\subsection{Architecture of Concurrent Servers}

The foundation of our design lies in the threading model introduced by
Anderson et al.~\cite{anderson1992}, which identified the core tension in
concurrent systems: kernel-level threads are expensive (context switches on
the order of 100\,$\mu$s) but integrate correctly with blocking I/O, while
user-level threads are an order of magnitude faster but stall the entire
process on any blocking system call. Our implementation resolves this tension
using POSIX \texttt{pthreads} with a fixed-size pool, accepting kernel-thread
cost in exchange for correct semantics when workers block on file I/O.

Building on this threading foundation, Welsh, Culler, and Brewer introduced
SEDA~\cite{welsh2001}---a Staged Event-Driven Architecture that partitions
server logic into a pipeline of stages, each with its own thread pool and
queue. SEDA's Haboob HTTP server implementation used a form of shortest
connection first scheduling within each stage, directly motivating our SFF
policy. The producer-consumer bounded buffer at the heart of our design is a
single-stage simplification of SEDA's multi-stage pipeline.

von Behren et al.\ later challenged SEDA's event-driven philosophy with
Capriccio~\cite{vonbehren2003}, arguing that user-level thread systems can
match event-driven performance while retaining the simpler sequential
programming model. Capriccio introduced resource-aware scheduling, where the
scheduler used compiler-annotated resource usage predictions to make smarter
thread selection decisions---conceptually analogous to our use of
\texttt{stat()} to obtain file size before enqueuing a request.

\subsection{Size-Based Scheduling Theory}

The theoretical backbone of our research extension is the work of
Harchol-Balter et al.~\cite{harchol2003}, who formalized size-based
scheduling for web servers as an application of
Shortest-Remaining-Processing-Time (SRPT) queuing theory. The key insight is
that, under heavy-tailed workloads, prioritizing short jobs (small files)
dramatically reduces mean sojourn time because the population of competing
small jobs is large. The paper's trace-driven simulations on real web server
logs showed mean response time improvements of 2--7$\times$ over FCFS at
moderate-to-high loads.

Schroeder and Harchol-Balter~\cite{schroeder2006} extended this analysis to
overload conditions, proving that FCFS response times explode exponentially as
load approaches server capacity, while SFF-style policies remain stable. This
is critical for our experimental design: we must measure not only average
conditions but also behavior as concurrent connections increase toward
saturation. The paper also characterized the ``starvation myth''---large
requests are rarely penalized significantly under SFF because they are
statistically infrequent in heavy-tailed workloads.

The heavy-tailed nature of web traffic itself was characterized by Crovella
and Bestavros~\cite{crovella1997}, who demonstrated through empirical analysis
of HTTP traces that file size distributions are Pareto-like with shape
parameter $\alpha \approx 1.1$. This self-similarity is the fundamental
property that makes size-based scheduling effective and is the reason our
benchmarking rig generates workloads with mixed small and large files rather
than uniform sizes.

\subsection{Modern Production Scheduling Systems}

Recent systems work validates and extends the scheduling ideas above at
production scale. Ousterhout et al.'s Shenango~\cite{ousterhout2019} achieves
microsecond-scale thread scheduling via a dedicated IOKernel that reallocates
CPU cores across applications as request rates change. While Shenango operates
at a finer granularity than our thread pool model, its architecture directly
parallels ours: the IOKernel is analogous to our main acceptor thread, and the
per-application runtime pools are analogous to our worker thread pool.
Shenango's evaluation confirmed that scheduling policy---not mechanism
overhead---is the primary determinant of tail latency.

Humphries et al.'s ghOSt~\cite{ghost2021}, deployed in production at Google,
decouples scheduling policy from the kernel enforcement mechanism entirely,
allowing user-space agents to implement FIFO, priority, or preemptive policies
by submitting transactions to the kernel. This validates our design choice of
making scheduling configurable via a runtime flag (\texttt{-s schedalg})
rather than hardcoding one policy. Demoulin et al.'s
Perséphone~\cite{persephone2021} demonstrated that size-based scheduling at
the microsecond scale---for workloads with heavy-tailed service time
distributions similar to our file size distribution---continues to provide
substantial tail latency improvements over FIFO, directly corroborating the
theoretical predictions of Harchol-Balter.

Lampson's foundational design heuristics~\cite{lampson1983} also informed our
architectural choices. The principle of splitting resources in a fixed way
justifies our fixed-size thread pool over dynamic thread creation, and the
advice to handle the normal and worst case separately underpins our SFF
rationale: the normal case (small files) should complete immediately, while
the exceptional case (large files) must not block them.

% ═════════════════════════════════════════════════════════════════════════════
\section{System Design}
\label{sec:design}
% ═════════════════════════════════════════════════════════════════════════════

\subsection{Architectural Overview}

Our system implements the G.5 OSTEP concurrent web server specification. The
server accepts connections on a TCP socket and distributes HTTP requests to a
fixed pool of worker threads via a shared bounded buffer, following the
producer-consumer pattern. Three configurable parameters govern runtime
behavior: \texttt{-t~N} sets the thread pool size (default: 1), \texttt{-b~N}
sets the buffer capacity (default: 1), and \texttt{-s~\{FCFS|SFF\}} selects
the scheduling policy.

\subsection{Thread Pool and Bounded Buffer}

The bounded buffer is implemented as a fixed-capacity circular array protected
by a single \texttt{pthread\_mutex\_t} and two condition variables:
\texttt{notFull} (signaled by consumers after dequeue) and \texttt{notEmpty}
(signaled by the producer after enqueue). A fixed-size array is used rather
than a dynamic queue to ensure memory-safe operation under concurrent access
without requiring heap allocation inside the critical section.

Each slot in the buffer holds a \texttt{request\_t} struct that carries all
the information a worker thread needs to service a request without accessing
any shared state after dequeue. The \texttt{client\_fd} field is the accepted
socket descriptor passed from the main thread; \texttt{filename} holds the
sanitized URI path extracted during initial parsing; \texttt{file\_size}
stores the result of the \texttt{stat()} call made by the producer before
acquiring the mutex; and \texttt{arrival\_ms} records the enqueue timestamp
used by FCFS to preserve arrival order. The struct is defined as follows:

\begin{lstlisting}
typedef struct {
    int    client_fd;     /* accepted socket descriptor */
    char   filename[256]; /* parsed URI path           */
    off_t  file_size;     /* from stat(), set on enqueue */
    long   arrival_ms;    /* timestamp for FCFS order  */
} request_t;
\end{lstlisting}

The \texttt{file\_size} field is populated by the producer via a
\texttt{stat()} call \emph{before} entering the critical section, so that no
filesystem operation occurs while the mutex is held.

% Figure 1 placed AFTER the code listing to prevent float collision
\begin{figure}[t]
  \centering
  \includegraphics[width=\columnwidth]{figs/bounded-buffer-activity.pdf}
  \caption{Activity diagram of the bounded buffer. The listener thread
    (producer) enqueues connection descriptors; worker threads (consumers)
    dequeue according to the active scheduling policy. Mutex and condition
    variables (\texttt{notFull}/\texttt{notEmpty}) prevent busy-waiting and
    enforce mutual exclusion.}
  \label{fig:activity}
\end{figure}

\subsection{Scheduling Policies}

\textbf{FCFS.} The consumer removes the element at the buffer's logical
head---a standard circular-queue dequeue in $O(1)$.

\textbf{SFF.} The consumer scans all \texttt{count} slots in the buffer,
identifies the entry with the minimum \texttt{file\_size}, removes it, and
shifts the remaining elements to close the gap. This is an $O(n)$ scan where
$n$ is the current buffer occupancy. Given typical buffer sizes (8--64 slots),
the overhead is negligible compared to the disk I/O that follows.

The scheduling selection is encapsulated in a single function pointer assigned
at startup from the \texttt{-s} argument, so the rest of the worker code is
policy-agnostic.

\subsection{HTTP Request--Response Lifecycle}

Figure~\ref{fig:sequence} details the five phases of each request. On
connection establishment (Step~1), the main thread completes the TCP
three-way handshake and obtains a connected socket descriptor. This descriptor
is enqueued into the bounded buffer (Step~2). A worker thread dequeues it
(Step~3), reads and parses the HTTP request line, performs a \texttt{stat()}
to verify the file exists and obtain its size, then maps and transmits the
file contents with an appropriate status line (Step~4). The connection is
closed after the response unless \texttt{Connection: keep-alive} is negotiated
(Step~5). The server returns \texttt{200 OK} for valid requests and
\texttt{404 Not Found} for missing resources.

\begin{figure}[t]
  \centering
  \includegraphics[width=\columnwidth]{figs/http-sequence.pdf}
  \caption{Sequence diagram of the full HTTP request--response lifecycle,
    including the TCP three-way handshake (Step~1), request transmission
    (Step~2), server processing (Step~3), response delivery (Step~4), and
    connection teardown (Step~5).}
  \label{fig:sequence}
\end{figure}

\subsection{Request-to-Job Dispatch}

Figure~\ref{fig:lifecycle} illustrates the system from a kernel-interaction
perspective. The Dispatcher component in our implementation is the bounded
buffer's dequeue logic. This is where our research intervention occurs: by
changing the selection criterion at Step~4 from arrival-order (FCFS) to
minimum-file-size (SFF), we hypothesize a significant reduction in mean
response time for small-file clients under load.

\begin{figure}[t]
  \centering
  \includegraphics[width=\columnwidth]{figs/request-job-lifecycle.pdf}
  \caption{The request-to-job lifecycle across the Client, Kernel, Server
    Process, Dispatcher, and Execution Unit. Our research intervention occurs
    at Step~4 (Job Dispatch), where the Dispatcher applies either FCFS or SFF
    to select the next request from the bounded buffer.}
  \label{fig:lifecycle}
\end{figure}

\subsection{Source Organization}

The codebase is organized into six focused modules, each with a single
responsibility, plus a top-level \texttt{Makefile}. This separation ensures
that the bounded buffer, thread pool, and HTTP logic remain independently
testable and that scheduling policy changes are confined to
\texttt{src/buffer.c} without touching any other file.
Table~\ref{tab:files} lists each module and its role.

\begin{table}[h]
\centering
\caption{Source file responsibilities}
\label{tab:files}
\small
\begin{tabular}{@{}ll@{}}
\toprule
\textbf{File} & \textbf{Responsibility} \\
\midrule
\texttt{main.c}              & Arg parsing, socket setup, accept loop \\
\texttt{bounded\_buffer.c/h} & Buffer, mutex, condvars, FCFS/SFF \\
\texttt{thread\_pool.c/h}    & Pool init, thread creation, teardown \\
\texttt{http.c/h}            & HTTP parsing, file send, status codes \\
\texttt{mock\_client.c}      & Correctness testing client \\
\texttt{Makefile}            & \texttt{make} / \texttt{make clean} \\
\bottomrule
\end{tabular}
\end{table}

\subsection{Concurrency Correctness}

All accesses to the shared bounded buffer pass through the mutex. Condition
variables are checked inside \texttt{while} loops (not \texttt{if} statements)
to guard against spurious wakeups. The producer holds the mutex only while
inserting into the array; the \texttt{stat()} call is performed beforehand.
Workers hold the mutex only while scanning/removing from the buffer; all file
I/O occurs after releasing the lock. This minimizes critical section duration
and avoids holding the mutex across blocking calls.

% ═════════════════════════════════════════════════════════════════════════════
\section{Methodology}
\label{sec:methodology}
% ═════════════════════════════════════════════════════════════════════════════

\subsection{Benchmarking Rig}

All performance measurements are collected externally, treating the server
as a black box. We use \texttt{wrk}~\cite{wrk} as our primary load generator
because it supports configurable concurrency, duration-based runs, and Lua
scripting for mixed-file-size workloads. Scripts are located in
\texttt{/benchmarks/} and are fully automated: a single invocation of
\texttt{benchmark.sh} creates the small and large test files using the
\texttt{dd} Unix utility, starts the server, runs all experiments, tears
down the server, and writes raw output to \texttt{/benchmarks/raw/}.

\subsection{Experimental Variables}

\textbf{Independent variables:}
\begin{itemize}
  \item \textbf{Scheduling policy:} FCFS vs.\ SFF (primary comparison).
  \item \textbf{Concurrency level:} number of simultaneous \texttt{wrk}
    connections: $\{1, 10, 25, 50, 100, 200\}$.
  \item \textbf{File size mix:} uniform small (4\,KB), uniform large (1\,MB),
    and heavy-tailed (90\% small / 10\% large by request count).
\end{itemize}

\textbf{Dependent variables:}
\begin{itemize}
  \item Throughput (requests per second).
  \item Response time: mean, p50, p95, p99 (milliseconds).
  \item Error rate (connection resets, timeouts).
\end{itemize}

\textbf{Controlled variables:} thread pool size fixed at 4; buffer capacity
fixed at 16; same VM and OS image; no background processes during measurement;
server and client run on the same host to eliminate network variability.

\subsection{Experimental Protocol}

Each configuration is run for 30 seconds. Before each individual trial's
measurement window, a 5-second warm-up period is performed against the
already-running server (since the server is restarted between trials,
each trial starts cold and requires its own warm-up). We collect five
independent trials per configuration, reporting the mean and standard
deviation. Between trials the server is restarted to clear any warm state.
Configurations where the server reaches connection saturation --- indicated
by a surge in connection refusal errors reported by \texttt{wrk} --- are
noted as such in the results rather than omitted, as saturation behavior
itself is a meaningful data point for comparing FCFS and SFF under overload.

\subsection{Environment}

\begin{table}[h]
\centering
\caption{Measurement environment}
\label{tab:env}
\begin{tabular}{@{}ll@{}}
\toprule
\textbf{Component} & \textbf{Specification} \\
\midrule
Host machine    & Windows 10 \\
Virtualization  & Oracle VirtualBox \\
Guest OS        & Ubuntu 22.04 LTS (Linux 5.15) \\
Guest resources & 4 vCPUs, 8\,GB RAM \\
Compiler        & GCC 11.4, \texttt{-O2} \\
Load generator  & \texttt{wrk} 4.1.03\\
\bottomrule
\end{tabular}
\end{table}

% ═════════════════════════════════════════════════════════════════════════════
\section{Preliminary Results}
\label{sec:results}
% ═════════════════════════════════════════════════════════════════════════════

\subsection{Baseline System Verification}

The base system (FCFS only) was verified to correctly serve static files to a
browser, handle concurrent connections without data corruption or deadlock,
and return appropriate HTTP status codes (\texttt{200}, \texttt{404}).

\subsection{Baseline Performance Numbers}

Table~\ref{tab:baseline} shows baseline FCFS measurements under the three
file-size workloads at selected concurrency levels. 

\begin{table}[h]
\centering
\caption{FCFS baseline: mean response time (ms) and throughput (req/s)}
\label{tab:baseline}
\begin{tabular}{@{}ccccc@{}}
\toprule
\textbf{Workload} & \textbf{Conc.} & \textbf{Mean RT (ms)} & \textbf{p99 (ms)} & \textbf{Req/s} \\
\midrule
Small (4\,KB)  &  10 & 0.56  & 3.15  & 15{,}122 \\
Small (4\,KB)  & 100 & 3.55  & 6.81  & 26{,}715 \\
Large (1\,MB)  &  10 & 6.64  & 17.97 & 1{,}462  \\
Large (1\,MB)  & 100 & 56.34 & 87.61 & 1{,}765  \\
Heavy-tailed   &  10 & 1.21  & 5.97  & 8{,}074  \\
Heavy-tailed   & 100 & 8.79  & 16.23 & 11{,}149 \\
\bottomrule
\end{tabular}
\end{table}

Values represent the mean across five independent trials. Raw
\texttt{wrk} output is committed to \texttt{/benchmarks/raw/}.

\subsection{Expected Trends}

Based on theoretical predictions from Harchol-Balter et al.~\cite{harchol2003}
and Schroeder and Harchol-Balter~\cite{schroeder2006}, we anticipate the
following patterns in the baseline data:

\begin{itemize}
  \item Under the \textbf{small-file} workload, mean response time will scale
    nearly linearly with concurrency and remain low, as all requests are
    uniformly cheap.
  \item Under the \textbf{large-file} workload, even moderate concurrency will
    cause mean response time to rise sharply, as each active worker ties up a
    thread for the full I/O duration.
  \item Under the \textbf{heavy-tailed} workload, mean response time will be
    dominated by the occasional large-file request blocking small requests
    behind it---the precise regime where SFF is predicted to help most.
\end{itemize}

Table~\ref{tab:expected} summarizes the hypothesized improvement of SFF over
FCFS for Sprint~4 comparison, drawn from~\cite{harchol2003}.

\begin{table}[h]
\centering
\caption{Hypothesized SFF improvement over FCFS (heavy-tailed workload)}
\label{tab:expected}
\begin{tabular}{@{}ccc@{}}
\toprule
\textbf{Load level} & \textbf{FCFS mean RT} & \textbf{Expected SFF gain} \\
\midrule
Low  ($<\!50\%$ capacity)  & Low       & Minimal ($<\!10\%$)        \\
Med  (${\sim}70\%$)        & Moderate  & ${\sim}1.5$--$2\times$     \\
High (${\sim}80\%$)        & High      & ${\sim}3$--$4\times$       \\
Near sat.\ ($>\!90\%$)     & Very high & $>\!6\times$               \\
\bottomrule
\end{tabular}
\end{table}

% ═════════════════════════════════════════════════════════════════════════════
\section{Conclusion}
\label{sec:conclusion}
% ═════════════════════════════════════════════════════════════════════════════

This paper has described the design and baseline measurement of a concurrent
POSIX web server implementing the G.5 OSTEP specification. Our system
correctly serves static files to concurrent clients using a fixed-size thread
pool, a producer-consumer bounded buffer, and configurable FCFS or SFF
scheduling. Three UML diagrams---the bounded buffer activity diagram
(Figure~\ref{fig:activity}), the HTTP request-response sequence
(Figure~\ref{fig:sequence}), and the request-to-job lifecycle
(Figure~\ref{fig:lifecycle})---document the critical paths through the system.

The related work survey established that size-based scheduling is theoretically
well-motivated (Harchol-Balter~\cite{harchol2003}; Schroeder~\cite{schroeder2006}),
architecturally principled (Welsh~\cite{welsh2001}; Lampson~\cite{lampson1983}),
and validated at production scale in modern systems
(ghOSt~\cite{ghost2021}; Perséphone~\cite{persephone2021};
Shenango~\cite{ousterhout2019}).

In Sprint~4, we will implement SFF scheduling alongside the FCFS baseline,
execute the full benchmarking protocol defined in
Section~\ref{sec:methodology}, populate Table~\ref{tab:baseline} with measured
data, and provide an evidence-based answer to our research question regarding
the practical impact of scheduling policy under heavy-tailed web workloads.

% ═════════════════════════════════════════════════════════════════════════════
\begin{thebibliography}{99}
% ═════════════════════════════════════════════════════════════════════════════

\bibitem{ostep}
R. H. Arpaci-Dusseau and A. C. Arpaci-Dusseau,
\textit{Operating Systems: Three Easy Pieces}, v1.10.
Arpaci-Dusseau Books, 2023.
\url{https://pages.cs.wisc.edu/~remzi/OSTEP/}

\bibitem{welsh2001}
M. Welsh, D. Culler, and E. Brewer,
``SEDA: An Architecture for Well-Conditioned, Scalable Internet Services,''
in \textit{Proc. 18th ACM Symp. Operating Systems Principles (SOSP~'01)},
Banff, Canada, Oct. 2001, pp.~230--243.

\bibitem{harchol2003}
M. Harchol-Balter, B. Schroeder, N. Bansal, and M. Agrawal,
``Size-Based Scheduling to Improve Web Performance,''
\textit{ACM Trans. Comput. Syst.}, vol.~21, no.~2, pp.~207--233, May 2003.

\bibitem{schroeder2006}
B. Schroeder and M. Harchol-Balter,
``Web Servers Under Overload: How Scheduling Can Help,''
\textit{ACM Trans. Internet Technol.}, vol.~6, no.~1, pp.~20--52, Feb. 2006.

\bibitem{crovella1997}
M. E. Crovella and A. Bestavros,
``Self-Similarity in World Wide Web Traffic: Evidence and Possible Causes,''
\textit{IEEE/ACM Trans. Networking}, vol.~5, no.~6, pp.~835--846, Dec. 1997.

\bibitem{anderson1992}
T. E. Anderson, B. N. Bershad, E. D. Lazowska, and H. M. Levy,
``Scheduler Activations: Effective Kernel Support for the User-Level
Management of Parallelism,''
\textit{ACM Trans. Comput. Syst.}, vol.~10, no.~1, pp.~53--79, Feb. 1992.

\bibitem{vonbehren2003}
R. von Behren, J. Condit, F. Zhou, G. C. Necula, and E. Brewer,
``Capriccio: Scalable Threads for Internet Services,''
in \textit{Proc. 19th ACM Symp. Operating Systems Principles (SOSP~'03)},
Bolton Landing, NY, Oct. 2003, pp.~268--281.

\bibitem{lampson1983}
B. W. Lampson,
``Hints for Computer System Design,''
\textit{ACM SIGOPS Oper. Syst. Rev.}, vol.~17, no.~5, pp.~33--48, Oct. 1983.

\bibitem{ghost2021}
J. T. Humphries \textit{et al.},
``ghOSt: Fast \& Flexible User-Space Delegation of Linux Scheduling,''
in \textit{Proc. 28th ACM Symp. Operating Systems Principles (SOSP~'21)},
Virtual, Oct. 2021, pp.~588--604.

\bibitem{persephone2021}
H. M. Demoulin \textit{et al.},
``When Idling Is Ideal: Optimizing Tail-Latency for Heavy-Tailed Datacenter
Workloads with Perséphone,''
in \textit{Proc. 28th ACM Symp. Operating Systems Principles (SOSP~'21)},
Virtual, Oct. 2021, pp.~621--637.

\bibitem{ousterhout2019}
A. Ousterhout, J. Fried, J. Behrens, A. Belay, and H. Balakrishnan,
``Shenango: Achieving High CPU Efficiency for Latency-sensitive Datacenter
Workloads,''
in \textit{Proc. 16th USENIX Symp. Networked Systems Design and
Implementation (NSDI~'19)},
Boston, MA, Feb. 2019, pp.~361--378.

\bibitem{wrk}
W. Glozer, ``wrk --- a HTTP benchmarking tool,'' GitHub, 2023.
\url{https://github.com/wg/wrk}

\end{thebibliography}

\end{document}